% =============================================================================
% Part 5: Discussion + Conclusion and Future Scope
% (Related Work has moved to Section 2. Broader Impact is absorbed into the
%  Discussion as a final ethics paragraph.)
% =============================================================================

\section{Discussion}
\label{sec:discussion}

\paragraph{Theoretical contribution:}
The cleanest empirical signal we measured---a recognition advantage on tagged facts---did not survive transformation into autoregressive generation, and did not survive a deliberate attempt to teach the model to make it survive. We interpret this as a structural property of attention-augmented language models: stored memory K/V provide attention-time logit perturbations that can shift a discrete classification but cannot consistently steer token-by-token sampling against a strong pretrained prior. The warm-up experiment (Section~\ref{sec:results_warmup}) turns this from an observation into a manipulation: training that demonstrably taught memory-conditioned prediction raised recognition while leaving recall at floor---on four models across three families---dissociating the two capacities experimentally rather than merely correlationally. Two further conceptual contributions follow. The self-grading result quantifies a failure mode of proxy validation at $93$--$100\%$ false confirmation on every model tested, which is severe enough to reframe surprise-reduction as unusable in isolation rather than merely imperfect. And the horizon programme shows that a headline comparison between a constrained and an unconstrained architecture is doubly unstable---it inverts with the number of cycles \emph{and} with the random seed, the constrained arm's damage plateau varying over three orders of magnitude on identical hyperparameters---which bears on how the continual-learning literature reports results generally.

\paragraph{Practical implications:}
For practitioners, four operational lessons follow. First, direct-write KV memory is usable for \emph{recognition}-style tasks on a frozen backbone---classification, multiple choice, similarity judgement---but not for free-form generation, on any family we tested, and a cheap retrieval-aware warm-up does not change that; closing the gap appears to require memory-augmented training at the scale of \citet{wu2022memorizing}. Second, any system that validates its own learning against a proxy should be audited against an external criterion before that proxy is trusted; on every model, nearly all candidates our proxy certified were rejected by a direct test. Third, continual-learning comparisons should be reported across both a horizon sweep and a seed spread: our nine ten-cycle run-pairs show the winner at a fixed horizon is decided as much by the seed lottery as by the architecture. Fourth, when adopting a compound safety mechanism, ablate it: at 7B scale the EWC-only baseline achieves a better recall-per-preservation ratio than the full pipeline (the effect weakens at 1.5B), which suggests the additional machinery is not carrying its weight in this regime.

\paragraph{Challenges:}
Several scope constraints bound the strength of our claims:

\begin{itemize}[leftmargin=2.4em,itemsep=0.15em,topsep=0.2em]
    \item[(i)] We test four models across three families (Qwen2.5-7B/1.5B, Llama-3.1-8B, Mistral-7B), which replicates the recall floor, the self-grading gap, and the warm-up null cross-family---but all are dense decoder-only architectures at $\leq$8B; mixture-of-experts and much larger scales remain untested.
    \item[(ii)] We test \emph{synthetic} facts in 10 templates; natural document content may differ.
    \item[(iii)] We test a \emph{single} LoRA configuration (rank 16, V/O projections, top-third layers); other adapter geometries may behave differently.
    \item[(iv)] We force \emph{eager attention}; SDPA and FlashAttention paths handle 4D additive masks differently and we did not characterize them after the gating mechanism was finalized.
    \item[(v)] The ten-cycle programme runs nine seeded pairs across the four models, which is enough to establish that the plateau level is high-variance but not enough to characterise its \emph{distribution}; a many-seed study of the plateau lottery is open.
    \item[(vi)] The warm-up is tested at \emph{one scale} ($\leq$300 steps, 400 passages) with a single, untuned learning rate; the Llama/Mistral collapses show this matters. A negative result at this budget does not bound what full memory-augmented pretraining would achieve.
\end{itemize}

Where variance was measured repeatedly it separates cleanly into two regimes. The single-cycle metrics are tight (DRA SD $\leq 0.003$ on every model; the proxy-vs-recall validation gap exceeds any seed fluctuation by two orders of magnitude), so the diagnostic findings are robust. The long-horizon metrics are the opposite---plateau levels vary by up to three orders of magnitude across seeds of one configuration---and we report that variance as a finding rather than averaging over it. The strongest claim our evidence supports is: on these models, on this dataset, with these safety hyperparameters, we observe the patterns reported.

\paragraph{Ethical considerations:}
Continual-learning systems for language models, if they work, would reduce retraining costs substantially and enable models to learn from interaction. They would also raise privacy and safety concerns: persistent memory of user inputs, model memorization of sensitive information, and asymmetric power between systems that can update themselves and users who cannot. Our specific architecture stores user inputs (tagged spans) in the model's KV bank and consolidates a subset into permanent weights; deployed naively, this would amount to silent on-device data retention. We argue that any deployed continual-learning system must include explicit user controls over what is encoded, what is consolidated, and what is permanent, alongside mechanisms for selective forgetting on user request. We note additionally that the self-grading result has a safety dimension: a system that certifies its own learning against an internal proxy may report successful consolidation of information it cannot actually reproduce, or fail to report retention of information it can.

\section{Conclusion and Future Scope}
\label{sec:conclusion}

We implemented and empirically evaluated a biologically-faithful memory consolidation architecture for transformer language models, then attempted to repair the central failure we found, replicating every headline result across four models from three families. Our experiments show that direct-write KV memory injection produces a measurable recognition signal in pretrained transformers but does not transfer to free-form recall---a floor of $\mathrm{DRA} \approx 0.006$ on every model, from models that follow the same facts in context at up to $0.99$; that a retrieval-aware warm-up phase, the repair the diagnosis implies, learns its objective and improves recognition without moving recall in any of sixteen runs; that internal surprise-based validation confirms consolidations at a $93$--$100\%$ false-confirmation rate against a direct recall test; that no single-cycle configuration achieves useful recall within an acceptable preservation budget; that over nine seeded ten-cycle run-pairs the safety machinery's damage-capping clip yields a universally-shaped but lottery-levelled preservation plateau, making fixed-horizon single-run comparisons unreliable; and that five specific assumptions in the original formalization required empirical correction before the architecture would function at scale. The key open problem is unchanged and now better specified: bridging from attention-based recognition into weights that support autoregressive generation, without sacrificing base-model preservation.

\paragraph{Future scope:}
Our results give follow-up work specific targets rather than vague directions:

\begin{itemize}[leftmargin=2.4em,itemsep=0.15em,topsep=0.2em]
    \item[(i)] \emph{Characterising the plateau lottery}: our nine run-pairs establish that the clip-induced preservation plateau has huge seed variance; a many-seed study could map its distribution, identify what in the optimisation path sets the level, and test whether early trajectory statistics predict it---which would turn the lottery into a controllable quantity.
    \item[(ii)] \emph{Retrieval-aware consolidation at scale}: our warm-up failed at $\leq$300 steps with a $36$-parameter gate and a rank-16 adapter, on four models. The informative next experiment is the same intervention at the training scale of Memorizing Transformers with per-family-tuned learning rates, which would establish whether the gap is closable in principle within a consolidation pipeline or requires pretraining-time memory augmentation.
    \item[(iii)] \emph{External validation criteria}: replace surprise-reduction with the direct recall gate inside the consolidation loop, and measure whether a system that cannot self-certify falsely also consolidates more effectively.
    \item[(iv)] \emph{Ablation-first safety design}: given that EWC alone outperformed the compound mechanism at 7B scale, rebuild the safety machinery additively, admitting each component only where it demonstrably improves the recall-per-preservation frontier.
    \item[(v)] \emph{Beyond dense $\leq$8B decoders}: extend the cross-family matrix to mixture-of-experts architectures and larger scales, where both the recognition signal and the plateau dynamics may differ.
\end{itemize}

We release the implementation, formalization, and full experimental logs to enable that work.
